% Options for packages loaded elsewhere
\PassOptionsToPackage{unicode}{hyperref}
\PassOptionsToPackage{hyphens}{url}
%
\documentclass[
  12pt,
  a4paper,
]{article}
\usepackage{amsmath,amssymb}
\usepackage{iftex}
\ifPDFTeX
  \usepackage[T1]{fontenc}
  \usepackage[utf8]{inputenc}
  \usepackage{textcomp} % provide euro and other symbols
\else % if luatex or xetex
  \usepackage{unicode-math} % this also loads fontspec
  \defaultfontfeatures{Scale=MatchLowercase}
  \defaultfontfeatures[\rmfamily]{Ligatures=TeX,Scale=1}
\fi
\usepackage{lmodern}
\ifPDFTeX\else
  % xetex/luatex font selection
\fi
% Use upquote if available, for straight quotes in verbatim environments
\IfFileExists{upquote.sty}{\usepackage{upquote}}{}
\IfFileExists{microtype.sty}{% use microtype if available
  \usepackage[]{microtype}
  \UseMicrotypeSet[protrusion]{basicmath} % disable protrusion for tt fonts
}{}
\makeatletter
\@ifundefined{KOMAClassName}{% if non-KOMA class
  \IfFileExists{parskip.sty}{%
    \usepackage{parskip}
  }{% else
    \setlength{\parindent}{0pt}
    \setlength{\parskip}{6pt plus 2pt minus 1pt}}
}{% if KOMA class
  \KOMAoptions{parskip=half}}
\makeatother
\usepackage{xcolor}
\usepackage[top=20mm, left=20mm, right=20mm, bottom=20mm]{geometry}
\usepackage{longtable,booktabs,array}
\usepackage{calc} % for calculating minipage widths
% Correct order of tables after \paragraph or \subparagraph
\usepackage{etoolbox}
\makeatletter
\patchcmd\longtable{\par}{\if@noskipsec\mbox{}\fi\par}{}{}
\makeatother
% Allow footnotes in longtable head/foot
\IfFileExists{footnotehyper.sty}{\usepackage{footnotehyper}}{\usepackage{footnote}}
\makesavenoteenv{longtable}
\usepackage{graphicx}
\makeatletter
\def\maxwidth{\ifdim\Gin@nat@width>\linewidth\linewidth\else\Gin@nat@width\fi}
\def\maxheight{\ifdim\Gin@nat@height>\textheight\textheight\else\Gin@nat@height\fi}
\makeatother
% Scale images if necessary, so that they will not overflow the page
% margins by default, and it is still possible to overwrite the defaults
% using explicit options in \includegraphics[width, height, ...]{}
\setkeys{Gin}{width=\maxwidth,height=\maxheight,keepaspectratio}
% Set default figure placement to htbp
\makeatletter
\def\fps@figure{htbp}
\makeatother
\setlength{\emergencystretch}{3em} % prevent overfull lines
\providecommand{\tightlist}{%
  \setlength{\itemsep}{0pt}\setlength{\parskip}{0pt}}
\setcounter{secnumdepth}{5}
% definitions for citeproc citations
\NewDocumentCommand\citeproctext{}{}
\NewDocumentCommand\citeproc{mm}{%
  \begingroup\def\citeproctext{#2}\cite{#1}\endgroup}
\makeatletter
 % allow citations to break across lines
 \let\@cite@ofmt\@firstofone
 % avoid brackets around text for \cite:
 \def\@biblabel#1{}
 \def\@cite#1#2{{#1\if@tempswa , #2\fi}}
\makeatother
\newlength{\cslhangindent}
\setlength{\cslhangindent}{1.5em}
\newlength{\csllabelwidth}
\setlength{\csllabelwidth}{3em}
\newenvironment{CSLReferences}[2] % #1 hanging-indent, #2 entry-spacing
 {\begin{list}{}{%
  \setlength{\itemindent}{0pt}
  \setlength{\leftmargin}{0pt}
  \setlength{\parsep}{0pt}
  % turn on hanging indent if param 1 is 1
  \ifodd #1
   \setlength{\leftmargin}{\cslhangindent}
   \setlength{\itemindent}{-1\cslhangindent}
  \fi
  % set entry spacing
  \setlength{\itemsep}{#2\baselineskip}}}
 {\end{list}}
\usepackage{calc}
\newcommand{\CSLBlock}[1]{\hfill\break\parbox[t]{\linewidth}{\strut\ignorespaces#1\strut}}
\newcommand{\CSLLeftMargin}[1]{\parbox[t]{\csllabelwidth}{\strut#1\strut}}
\newcommand{\CSLRightInline}[1]{\parbox[t]{\linewidth - \csllabelwidth}{\strut#1\strut}}
\newcommand{\CSLIndent}[1]{\hspace{\cslhangindent}#1}
\usepackage{pdflscape}
\usepackage{float}
\ifLuaTeX
  \usepackage{selnolig}  % disable illegal ligatures
\fi
\usepackage{bookmark}
\IfFileExists{xurl.sty}{\usepackage{xurl}}{} % add URL line breaks if available
\urlstyle{same}
\hypersetup{
  pdftitle={Behavioural Cost and COVID-Safe Behaviour Adherence Over Time in Australia},
  hidelinks,
  pdfcreator={LaTeX via pandoc}}

\title{Behavioural Cost and COVID-Safe Behaviour Adherence Over Time in Australia}
\usepackage{etoolbox}
\makeatletter
\providecommand{\subtitle}[1]{% add subtitle to \maketitle
  \apptocmd{\@title}{\par {\large #1 \par}}{}{}
}
\makeatother
\subtitle{Final Report}
\author{Indrasena Reddy Muthyala\\
\vspace{0.3cm}
Project supervisor: Anthony Mays}
\date{2026-08-04}

\begin{document}
\maketitle

\section{Declaration}\label{declaration}

In submitting this work I am indicating that I have read the
University's Academic Integrity Policy. I declare that all material in
this assessment is my own work except where there is clear
acknowledgement and reference to the work of others. I give permission
for this work to be reproduced and submitted to other academic staff for
educational purposes.

\section{Abstract}\label{abstract}

This study analysed how adherence to 14 COVID-safe behaviours changed
over time in Australia and whether behaviours involving greater
practical and social burden were less stable. Data was obtained from
the Australian component of the Imperial College London--YouGov COVID-19
Behaviour Tracker, containing 53,833 responses across 54 survey waves.
The behaviours were classified into low-cost, medium-cost, and high-cost groups.
Weighted adherence percentages were calculated across nine six-wave
blocks, and logistic regression models were fitted for five behaviours
showing the largest increases between qweeks 31--36 and qweeks 37--42.

Low-cost hygiene behaviours stayed high and stable,
while medium- and high-cost behaviours showed greater variation over
time. Mask wearing recorded the largest increase between the selected
periods, followed by avoiding guests and avoiding small, medium-sized,
and large gatherings. For all five modelled behaviours, including
state-by-period interaction terms significantly improved model fit,
showing that the size of the change in behaviour adherence differed across states and
territories.

The timing of these changes were consistent with the 2021 Delta outbreak
and stronger public-health measures in several states. The
findings show that adherence patterns varied with both behavioural
burden and state-specific policies.

\section{Introduction}\label{introduction}

Public-health measures depend partially on whether people adapt and
maintain the recommended protective behaviours (\citeproc{ref-ryan2025}{Ryan et al. 2025, 1--2}).
During the COVID-19 pandemic, Australians were asked to perform
behaviours ranging from routine hygiene practices to significant changes
in movement and social contact. These behaviours differed in the effort,
inconvenience, and disruption they caused, which may have influenced
how consistently they were followed (\citeproc{ref-hinssen2023}{Hinssen and Dohle 2023, secs. 1.1--1.2}).

Understanding these differences is important for future public-health
planning. Australian research has shown that hygiene behaviours can
remain relatively stable while behaviours involving physical
distancing change more quickly over time (\citeproc{ref-ayre2021}{Ayre et al. 2021, 1}).
Comparing multiple behaviours can therefore help identify which
protective actions remain stable and which respond more strongly to
changing outbreak and public-health conditions.

In this report, low-cost behaviours refer to actions requiring
little effort and disruption to daily life, while higher-cost behaviours
require greater changes to movement, social contact, or normal daily
activities.

The following research question is addressed:

\textbf{How did adherence to COVID-safe behaviours change over time in
Australia, and were higher-cost behaviours less stable than lower-cost
behaviours?}

The project also explores whether the selected change between qweeks
31--36 and qweeks 37--42 differed across Australian states and
territories. This report compares 14 protective behaviours
grouped by behavioural cost and to distinguish broad national
patterns from state-specific differences.

\section{Background}\label{background}

\subsection{Protective behaviours and adherence}\label{protective-behaviours-and-adherence}

Protective health behaviours reduce the risk of acquiring or
transmitting infection. During COVID-19, these included personal hygiene
behaviours such as washing hands with soap and water and covering the
nose and mouth when coughing or sneezing, as well as mask wearing,
avoiding close contact and gatherings, and limiting movement. Their
value depends on both effectiveness and continued public adherence.

Adherence can vary between behaviours and across time. Personal hygiene
behaviours can generally be performed within an individual's normal
routine, while behaviours such as avoiding gatherings or remaining at
home require larger changes to social and daily activities. Previous
research has reported associations between mask wearing, social
distancing, and hand hygiene, indicating that protective behaviours may
occur together (\citeproc{ref-ryan2025}{Ryan et al. 2025, 2}).

Research has identified several influences on protective behaviour,
including demographic characteristics, perceived illness threat, trust
in government, social norms, and public-health policy (\citeproc{ref-ryan2025}{Ryan et al. 2025, 1--2}). The importance of these factors may also change during a
pandemic. For example, before a mask mandate, mask wearing may primarily
reflect voluntary choice and perceived infection risk; after a mandate,
it may also reflect compliance with a legal requirement. Examining
several protective behaviours over time can therefore show whether
adherence follows a common pattern or changes according to the
particular action being considered.

\subsection{Behavioural cost and changes over time}\label{behavioural-cost-and-changes-over-time}

Protective behaviours differ in the effort and disruption required to
perform them. Some actions, such as washing hands or covering a cough,
can be added to an existing routine with relatively little change. Other
behaviours, such as avoiding gatherings, limiting travel, or remaining
at home, can restrict social contact, employment, movement, and access
to normal activities. These differences can be described as behavioural
cost or burden.

Protection Motivation Theory describes these barriers as \emph{response
costs}. Response costs can include inconvenience, discomfort, financial
expense, or disruption to everyday life. Higher expected response costs
may make protective behaviours more difficult to perform even when an
individual intends to follow the recommended action (\citeproc{ref-hinssen2023}{Hinssen and Dohle 2023, secs. 1.1--1.2}). In a longitudinal COVID-19 study, perceived costs were
negatively associated with subsequently reported physical distancing and
hand hygiene, although the corresponding direct effect was not
significant for mask wearing (\citeproc{ref-hinssen2023}{Hinssen and Dohle 2023, sec. 4}).

Behavioural burden may also help explain why protective behaviours
follow different patterns over time. Australian research conducted
during 2020 found that physical-distancing behaviours declined across
successive surveys, while hygiene behaviours remained comparatively
stable (\citeproc{ref-ayre2021}{Ayre et al. 2021, 1}). This distinction suggests that behaviours
requiring continued changes to movement and social interaction may be
less stable than behaviours that can become part of a daily routine.

Behavioural cost was considered one of several influences on adherence.
Perceived infection risk, confidence in performing the behaviour, social
expectations, and public-health rules may increase adherence despite the
burden involved. This report therefore examines both broad cost groups
and individual behaviours.

\subsection{Risk, policy, and the Australian context}\label{risk-policy-and-the-australian-context}

Adherence to protective behaviour can change when people notice a
greater risk of infection and when public-health rules become more
restrictive. Risk perception, trust in authorities, social expectations,
and government policy have all been identified as influences on COVID-19
protective behaviour (\citeproc{ref-ryan2025}{Ryan et al. 2025, 1--2}). These factors may affect
both whether a behaviour is adopted and whether it is maintained over
time.

Public-health rules can also change the practical meaning of a
behaviour. For example, mask wearing may shift from a voluntary
precaution to a legal requirement, while restrictions on gatherings,
household visitors, or movement may directly limit opportunities for
social contact. Changes in adherence may therefore reflect a combination
of voluntary responses to perceived risk and compliance with formal
rules.

In Australia, COVID-19 outbreaks and public-health measures differed
across states and territories. During the 2021 Delta period, New South
Wales, Victoria, and the Australian Capital Territory strengthened
stay-at-home orders, movement limits, mandatory mask requirements, and
restrictions on gatherings. Western Australia provided a contrasting
situation, returning towards pre-lockdown settings in July 2021 with
mask and crowd restrictions removed (\citeproc{ref-abs_delta_2022}{Australian Bureau of Statistics 2022}, section ``Delta''; \citeproc{ref-nsw_health_2021}{NSW Health 2021, paras. 1--3}; \citeproc{ref-vic_premier_2021}{State Government of Victoria 2021, paras. 1--7}; \citeproc{ref-act_gov_2021}{Barr and Stephen-Smith 2021, paras. 1--6}; \citeproc{ref-abc_wa_2021}{Perpitch 2021, paras. 1--5}).

Because these conditions varied between jurisdictions, a national
adherence estimate may combine different state-level patterns. This
provides the basis for examining both overall Australian trends and
differences between states and territories. However, temporal alignment
between policy changes and adherence cannot by itself establish that the
policies caused the observed behaviour.

\subsection{Logistic regression and model comparison}\label{logistic-regression-and-model-comparison}

Logistic regression is used to model an outcome with two possible
values, such as adherence (\texttt{1}) and non-adherence (\texttt{0}) (\citeproc{ref-hosmer2013}{Hosmer, Lemeshow, and Sturdivant 2013}).
For respondent \(i\), let

\[
p_i = P(Y_i = 1 \mid \mathbf{x}_i)
\]

represent the probability of adherence given a set of explanatory
variables \(\mathbf{x}_i\). Because a probability is restricted to values
between 0 and 1, logistic regression models the log-odds of the outcome:

\[
\operatorname{logit}(p_i)
=
\log\left(\frac{p_i}{1-p_i}\right)
=
\beta_0
+
\beta_1x_{i1}
+
\cdots
+
\beta_kx_{ik}.
\]

The quantity \(p_i/(1-p_i)\) is the odds of adherence. A regression
coefficient \(\beta_j\) represents the change in log-odds associated with
a one-unit change in \(x_j\), while the other variables in the model are
held constant. Exponentiating the coefficient gives an odds ratio:

\[
OR_j = \exp(\beta_j).
\]

An odds ratio greater than 1 represents higher odds of adherence, an
odds ratio below 1 represents lower odds, and an odds ratio equal to 1
represents equal odds. A coefficient p-value tests the null hypothesis
that \(\beta_j = 0\). A 95\% confidence interval describes the range of
coefficient values that are reasonably compatible with the fitted model.
Confidence intervals for odds ratios are obtained by exponentiating the
corresponding coefficient confidence limits.

An interaction term is used when the association between one explanatory
variable and the outcome may differ according to another variable. For
two variables \(X\) and \(Z\), a model containing an interaction can be
written as

\[
\operatorname{logit}(p_i)
=
\beta_0
+
\beta_1X_i
+
\beta_2Z_i
+
\beta_3(X_iZ_i).
\]

The coefficient \(\beta_3\) measures how the association between \(X\) and
the outcome differs across values or categories of \(Z\). In this study,
the interaction between survey period and state or territory was used to
assess whether the period change in adherence differed across
jurisdictions.

Logistic-regression coefficients are commonly estimated by maximum
likelihood (\citeproc{ref-hosmer2013}{Hosmer, Lemeshow, and Sturdivant 2013}). For independent binary observations, the
likelihood is

\[
L(\boldsymbol{\beta})
=
\prod_{i=1}^{n}
p_i^{y_i}
(1-p_i)^{1-y_i},
\]

where \(y_i\) is the observed binary outcome. The corresponding
log-likelihood is

\[
\ell(\boldsymbol{\beta})
=
\sum_{i=1}^{n}
\left[
y_i\log(p_i)
+
(1-y_i)\log(1-p_i)
\right].
\]

Maximum-likelihood estimation selects the coefficient values that
maximise the log-likelihood. For models fitted to the same observations,
a larger maximised log-likelihood indicates that the model assigns
greater probability to the observed outcomes.

A likelihood-ratio test compares two nested models, where the reduced
model is contained within the larger or full model (\citeproc{ref-hosmer2013}{Hosmer, Lemeshow, and Sturdivant 2013}). The
likelihood-ratio statistic is

\[
LR
=
2\left(
\ell_{\text{full}}
-
\ell_{\text{reduced}}
\right).
\]

Under the null hypothesis that the additional parameters do not improve
the model, the statistic is approximately chi-squared distributed. The
degrees of freedom equal the number of additional parameters in the full
model. A small p-value provides evidence that including those additional
parameters improves model fit.

\section{Methods}\label{methods}

\subsection{Overview of the analytical workflow}\label{overview-of-the-analytical-workflow}

The analysis was carried out in a sequence of linked steps designed to
answer the research question of how adherence to COVID-safe behaviours
changed over time in Australia and whether higher-cost behaviours were
less stable than lower-cost behaviours.

First, the Australian Behaviour Tracker data were cleaned and checked to
confirm the dataset structure and the availability of the variables
required for the analysis. Second, 14 protective-behaviour items were
selected and recoded into separate binary adherence outcomes. Third, the
behaviours were classified into low-, medium-, and high-cost groups
using a study-specific behavioural-burden framework.

Fourth, weighted descriptive adherence estimates were calculated for
individual behaviours and cost groups across nine six-wave blocks in
order to identify broad temporal patterns. Fifth, the descriptive
results were used to select qweeks 31--36 and qweeks 37--42 as the two
adjacent periods showing the clearest simultaneous increase in medium-
and high-cost adherence. Sixth, state-level descriptive comparisons were
used to examine whether the change between these two periods was similar
across jurisdictions.

Finally, logistic regression models were fitted for the five behaviours
showing the largest descriptive increases, and likelihood-ratio tests
were used to assess whether allowing the period change to differ across
states and territories improved model fit. Together, these steps
combined descriptive and regression-based analyses to examine both
overall adherence patterns and jurisdiction-specific differences.

\begin{center}

\textbf{Analytical workflow}

\vspace{0.2cm}

\fbox{
\parbox{0.9\textwidth}{
\centering
Data cleaning and structure checks \\[0.2cm]
$\downarrow$ \\[0.2cm]
Selection of 14 behaviour variables and binary adherence coding \\[0.2cm]
$\downarrow$ \\[0.2cm]
Behavioural-cost grouping \\[0.2cm]
$\downarrow$ \\[0.2cm]
Weighted descriptive analysis across nine six-wave blocks \\[0.2cm]
$\downarrow$ \\[0.2cm]
Selection of qweeks 31--36 and qweeks 37--42 \\[0.2cm]
$\downarrow$ \\[0.2cm]
State-level descriptive comparison \\[0.2cm]
$\downarrow$ \\[0.2cm]
Logistic regression and likelihood-ratio model comparison
}
}

\end{center}

\subsection{Data source and study design}\label{data-source-and-study-design}

This study used data from the Australian component of the Imperial
College London--YouGov COVID-19 Behaviour Tracker. The Behaviour Tracker
collected repeated cross-sectional survey data on protective behaviours,
attitudes, demographics, health, and responses to COVID-19 across
multiple countries between 2020 and 2022 (\citeproc{ref-jones2020}{Jones, Imperial College London Big Data Analytical Unit, and YouGov Plc 2020}; \citeproc{ref-ryan2025}{Ryan et al. 2025, 3}).

The working dataset contained 53,833 survey responses collected across
54 survey waves. Each row represented one survey response. All 53,833
responses had a unique RecordNo, with no duplicate values found in the
Australian dataset.

The survey was conducted repeatedly over time; however, the dataset did
not contain a stable person-level identifier that could be used to link
the same respondent across survey waves. The data were analysed as
repeated cross-sectional observations. Changes in adherence over time
therefore represent differences between the surveyed samples across
periods because individual respondents could not be followed across
waves.

Variables used in the final analysis included survey wave, survey end
date, state or territory, survey weight, and 14 COVID-safe behaviour
variables. Demographic variables were retained in the cleaned dataset
but were not included in the final regression models.

\subsection{Behaviour variables}\label{behaviour-variables}

The analysis included 14 protective-behaviour items from the
\texttt{i12\_health} series of the Behaviour Tracker questionnaire. Each item
measured how frequently the respondent reported performing a particular
protective behaviour. The original questionnaire wording and variable
identifiers were retained where possible to make the selection of
behaviours reproducible (\citeproc{ref-jones2020}{Jones, Imperial College London Big Data Analytical Unit, and YouGov Plc 2020}, codebook entries
\texttt{i12\_health\_1}--\texttt{i12\_health\_8} and \texttt{i12\_health\_11}--\texttt{i12\_health\_16}; \citeproc{ref-ryan2025}{Ryan et al. 2025, 4}).

The 14 protective-behaviour items were:

\begin{itemize}
\tightlist
\item
  wearing a mask outside the home (\texttt{i12\_health\_1});
\item
  washing hands (\texttt{i12\_health\_2});
\item
  using hand sanitiser (\texttt{i12\_health\_3});
\item
  covering coughs or sneezes (\texttt{i12\_health\_4});
\item
  avoiding potentially infected people (\texttt{i12\_health\_5});
\item
  avoiding going out (\texttt{i12\_health\_6});
\item
  avoiding healthcare settings (\texttt{i12\_health\_7});
\item
  avoiding public transport (\texttt{i12\_health\_8});
\item
  avoiding guests at home (\texttt{i12\_health\_11});
\item
  avoiding small gatherings (\texttt{i12\_health\_12});
\item
  avoiding medium-sized gatherings (\texttt{i12\_health\_13});
\item
  avoiding large gatherings (\texttt{i12\_health\_14});
\item
  avoiding crowded areas (\texttt{i12\_health\_15}); and
\item
  avoiding shops (\texttt{i12\_health\_16}).
\end{itemize}

Each item represented a distinct protective action, so a separate binary
adherence outcome was created for each behaviour.

\subsection{Binary adherence coding}\label{binary-adherence-coding}

All 14 behaviours were measured using five ordered response categories:
Always, Frequently, Sometimes, Rarely, and Not at all. In the original
Behaviour Tracker coding, these responses corresponded to values from 5
for Always to 1 for Not at all (\citeproc{ref-jones2020}{Jones, Imperial College London Big Data Analytical Unit, and YouGov Plc 2020}; \citeproc{ref-ryan2025}{Ryan et al. 2025, 4}).

For this study, each item was dichotomised separately. Always (\texttt{5}) and
Frequently (\texttt{4}) responses were classified as adherent (\texttt{1}). Sometimes
(\texttt{3}), Rarely (\texttt{2}), and Not at all (\texttt{1}) responses were classified as
not adherent (\texttt{0}).

Missing responses were retained as missing and were not assigned to
either adherence category. Consequently, the number of observations
available for analysis could differ between behaviours.

The same rule was applied separately to all 14 behaviour variables. All
14 binary outcomes were used in the descriptive analysis. The five
behaviours selected for the regression analysis were modelled separately
as binary dependent variables. The data checks, variable-selection
steps, and binary recoding procedure are documented in the Appendix
section titled \emph{Data preparation and code references}.

\subsection{Behavioural cost grouping}\label{behavioural-cost-grouping}

The 14 behaviours were classified into low-, medium-, and high-cost
groups using an exploratory behavioural-burden framework developed
specifically for this study. The framework has not been externally
validated. Behavioural cost was defined as the level of difficulty or
disruption involved in maintaining a protective behaviour during normal
daily life.

The five scoring dimensions operationalised the response-cost concept
discussed in the Background. Each behaviour was assessed according to
practical effort, social disruption, lifestyle restriction, emotional
burden, and financial burden. These dimensions represented the effort
required and the extent to which a behaviour could affect social
contact, daily activities, wellbeing, or finances.

The dimension scores were assigned by the researcher using these
definitions and the expected burden imposed by each behaviour. They were
not taken from a previously validated behavioural-cost scale.

Each dimension was assigned a score from 1 to 3, where 1 represented low
burden, 2 represented moderate burden, and 3 represented high burden.
The five scores were added to produce a total behavioural-cost score
ranging from 5 to 15. Behaviours with total scores of 5--6 were
classified as low cost, scores of 7--9 as medium cost, and scores of
10--15 as high cost.

Table \ref{tab:cost-groups} presents the final classification of the 14
behaviours. The complete scoring matrix, including the five individual
dimension scores and the total score for each behaviour, is provided in
Appendix Table \ref{tab:cost-scoring}.

\begin{longtable}[]{@{}ll@{}}
\caption{\label{tab:cost-groups} Classification of the 14 protective behaviours by
behavioural-cost group.}\tabularnewline
\toprule\noalign{}
Behaviour & Behavioural-cost group \\
\midrule\noalign{}
\endfirsthead
\toprule\noalign{}
Behaviour & Behavioural-cost group \\
\midrule\noalign{}
\endhead
\bottomrule\noalign{}
\endlastfoot
Covered coughs or sneezes & Low \\
Washed hands & Low \\
Used hand sanitiser & Low \\
Wore a mask outside the home & Medium \\
Avoided people with COVID-like symptoms & Medium \\
Avoided crowded areas & Medium \\
Avoided public transport & Medium \\
Avoided having guests at home & Medium \\
Avoided small gatherings & High \\
Avoided medium gatherings & High \\
Avoided large gatherings & High \\
Avoided going out & High \\
Avoided shops & High \\
Avoided healthcare settings & High \\
\end{longtable}

The classification was based solely on the five burden dimensions and
was completed independently of the observed adherence percentages. This
preserved a separation between the classification criteria and the
adherence patterns examined in the analysis.

The scoring framework, behaviour-specific rationale, implementation
details, and complete numerical scoring matrix are provided in Appendix
Section \emph{Behavioural-cost scoring matrix}.

\subsection{Survey timing and six-wave blocks}\label{survey-timing-and-six-wave-blocks}

The Behaviour Tracker dataset identified each survey wave using the
variable \texttt{qweek}. Each \texttt{qweek} value represents one survey wave. The
intervals between survey waves were not fixed at one calendar week. The
54 \texttt{qweek} values were converted into a numeric variable ranging from 1
to 54 so that the survey waves could be placed in their correct
chronological order.

The survey end-date variable was used to map each \texttt{qweek} to the
corresponding calendar period. This allowed the survey-based time
periods to be compared with the timing of COVID-19 outbreaks and
public-health measures in Australia. A complete mapping between the 54
survey waves and their calendar dates is provided in Appendix Table
\ref{tab:qweek-dates}.

For the descriptive analysis, the 54 survey waves were divided into nine
consecutive blocks containing six \texttt{qweek} values each: qweeks 1--6,
7--12, 13--18, 19--24, 25--30, 31--36, 37--42, 43--48, and 49--54. These
groups are referred to as six-wave blocks because each contains six
survey waves and may cover a different calendar duration.

The blocks reduced visual clutter while retaining equally sized groups
of survey waves. They do not cover equal calendar durations because data
collection did not occur at fixed weekly intervals.

After reviewing the nine-block trends, qweeks 31--36 and 37--42 were
selected because this adjacent pair showed the clearest simultaneous
increase in medium- and high-cost adherence. The period selection was
exploratory and was based on the observed descriptive pattern.

Based on the survey end dates, qweeks 31--36 covered 5 May to 27 July
2021, while qweeks 37--42 covered 28 July to 17 October 2021. These two
periods were used in the behaviour-level comparisons, state-level
comparisons, and logistic regression models. They are referred to
throughout the report as the earlier period and the later period,
respectively.

The survey-wave mapping, code references, and saved numerical output are
provided in Appendix Section \emph{Survey-wave and date mapping}.

\subsection{Weighted descriptive adherence estimates}\label{weighted-descriptive-adherence-estimates}

Weighted adherence percentages were calculated using the
respondent-level survey-weight variable supplied with the Behaviour
Tracker data (\citeproc{ref-jones2020}{Jones, Imperial College London Big Data Analytical Unit, and YouGov Plc 2020}). According to the repository documentation,
the supplied weighting variable was typically based on respondents' age,
gender, and region. The weights were used as provided in the dataset and
were not constructed or recalibrated as part of this study. Their
purpose in this analysis was to determine each respondent's contribution
to the weighted adherence estimates.

For behaviour \(b\), the weighted adherence proportion was calculated as

\[
\widehat{p}_{b}
=
\frac{\sum_{i=1}^{n} w_iY_{ib}}
     {\sum_{i=1}^{n} w_i},
\]

where \(Y_{ib}\) was the binary adherence outcome for respondent \(i\) and
behaviour \(b\), and \(w_i\) was the respondent's survey weight. Because
adherence was coded as \texttt{1} and non-adherence as \texttt{0}, the weighted mean
represented the weighted proportion classified as adherent. This
proportion was multiplied by 100 and reported as a percentage. A
respondent with a larger survey weight contributed proportionally more
to both the numerator and denominator of the estimate. The analytical
dataset did not contain the complete survey-design information required
to reconstruct the supplied weights independently.

The calculation was performed separately for each of the 14 behaviours
across the complete survey period and within each of the nine six-wave
blocks. For each estimate, observations were included only when both the
relevant adherence outcome and survey weight were available. All 53,833
responses had non-missing, positive survey weights. The handwashing and
avoiding-guests variables each contained 1,006 missing responses,
involving the same records, while the other 12 behaviour variables had
no missing responses. Missing behaviour responses were excluded only
from calculations involving that behaviour.

For the descriptive cost-group trends, each respondent's available
binary outcomes were averaged within the low-, medium-, and high-cost
groups. These respondent-level group rates were then summarised using
the survey weights within each six-wave block. The grouped estimates
were used only to describe broad patterns; individual behaviours
remained separate outcomes in the regression analysis.

The same weighted calculation was used to estimate behaviour-specific
adherence within each state or territory during qweeks 31--36 and qweeks
37--42. Changes between these periods were calculated in percentage
points by subtracting the earlier-period estimate from the later-period
estimate.

The corresponding code references and complete descriptive outputs are
provided in Appendix Section \emph{Complete descriptive estimates}.

\subsection{State-level descriptive comparison}\label{state-level-descriptive-comparison}

A state-level descriptive comparison was conducted to examine whether
the change observed in the national adherence estimates was similar
across Australian states and territories. This comparison was included
because the timing and extent of COVID-19 outbreaks and public-health
measures varied between jurisdictions.

The analysis compared qweeks 31--36 with qweeks 37--42. Based on the
survey end dates, these periods covered 5 May to 27 July 2021 and 28
July to 17 October 2021, respectively. For each state or territory
represented in the dataset, weighted adherence percentages were
calculated separately for the earlier and later periods using the same
weighting procedure described above.

For each jurisdiction and period, the respondent-level cost-group rates
described above were summarised using survey weights. Change was
calculated in percentage points as

\[
\text{Change}
=
\text{later-period adherence}
-
\text{earlier-period adherence}.
\]

A positive value represented an increase in weighted adherence between
the two periods, while a negative value represented a decrease.

This descriptive comparison summarised the direction and magnitude of
changes across jurisdictions. The regression analysis then examined
whether the period change differed across jurisdictions using
state-by-period interaction terms.

The unweighted sample sizes, weighted estimates, and percentage-point
changes for each jurisdiction are reported in Tables
\ref{tab:state-estimates} and \ref{tab:state-changes}.

\subsection{Logistic regression models}\label{logistic-regression-models}

The regression analysis examined whether adherence changed between the
two selected periods and whether the size of this change differed across
states and territories. Five behaviours were modelled separately because
they recorded the largest descriptive increases between qweeks 31--36
and qweeks 37--42:

\begin{itemize}
\tightlist
\item
  wearing a face mask outside the home;
\item
  avoiding having guests at home;
\item
  avoiding small social gatherings;
\item
  avoiding medium-sized social gatherings; and
\item
  avoiding large social gatherings.
\end{itemize}

For each behaviour, the outcome \(Y_i\) was coded as \texttt{1} for adherence and
\texttt{0} for non-adherence. The analysis used responses from qweeks 31--42.
The period variable was coded as \texttt{0} for qweeks 31--36 and \texttt{1} for
qweeks 37--42.

Two state-based models were fitted separately for each behaviour. Model
2 included survey period and state or territory as main effects:

\[
\operatorname{logit}\{P(Y_i = 1)\}
=
\beta_0
+
\beta_1\text{Period}_i
+
\boldsymbol{\beta}_2\text{State}_i.
\]

Model 3 included the same main effects together with state-by-period
interaction terms:

\[
\operatorname{logit}\{P(Y_i = 1)\}
=
\beta_0
+
\beta_1\text{Period}_i
+
\boldsymbol{\beta}_2\text{State}_i
+
\boldsymbol{\beta}_3
(\text{Period}_i \times \text{State}_i).
\]

The state variable was represented by indicator variables for the
Australian states and territories. New South Wales was the reference
jurisdiction, and qweeks 31--36 were the reference period. In Model 3,
\(\beta_1\) estimated the period change for New South Wales. Each
state-by-period interaction coefficient estimated how the period change
for another jurisdiction differed from the New South Wales change.

The models included 12,025 complete responses: 5,995 from qweeks 31--36
and 6,030 from qweeks 37--42. The five adherence outcomes, state,
period, and survey weight were complete for these responses.

The models were fitted in Python using a binomial generalised linear
model with a logit link. The respondent-level survey weights were
incorporated into model fitting. HC1 robust covariance estimates were
used to calculate coefficient standard errors, 95\% confidence intervals,
and coefficient p-values (\citeproc{ref-seabold2010}{Seabold and Perktold 2010}). The corresponding software
functions and implementation details are provided in Appendix Section
\emph{Complete Model 3 coefficient outputs}.

\subsection{Model comparison and statistical reporting}\label{model-comparison-and-statistical-reporting}

For each of the five behaviours, Model 2 was treated as the reduced
model and Model 3 as the full model. The models were compared using a
likelihood-ratio test to assess whether adding the state-by-period
interaction terms improved model fit.

Model 3 contained seven additional interaction parameters, corresponding
to the seven jurisdictions compared with New South Wales. Each
likelihood-ratio test therefore used seven degrees of freedom. Test
p-values were obtained from the chi-squared distribution
(\citeproc{ref-virtanen2020}{Virtanen et al. 2020}). The computational implementation is documented in
Appendix Section \emph{Likelihood-ratio test output}.

Model 3 coefficients, odds ratios, 95\% confidence intervals, and
coefficient p-values were extracted for statistical reporting. The
period coefficient for New South Wales and the state-by-period
interaction coefficients were used to describe the estimated period
differences across jurisdictions. Complete Model 3 coefficient outputs
are reported in Appendix Tables
\ref{tab:model3-mask-appendix}--\ref{tab:model3-large-appendix}, and
the likelihood-ratio test results are reported in Appendix Table
\ref{tab:lr-tests-appendix}.

\section{Results}\label{results}

\subsection{Overall weighted adherence by behaviour}\label{overall-weighted-adherence-by-behaviour}

\begin{figure}

{\centering \includegraphics[width=1\linewidth]{../graphs and tables/01_overall_adherence_by_behaviour} 

}

\caption{Overall weighted adherence percentage by COVID-safe behaviour.}\label{fig:overall-adherence-figure}
\end{figure}

Figure 1 presents the overall weighted adherence percentages for the 14
protective behaviours. The highest adherence was recorded for covering
coughs or sneezes (90.9\%) and washing hands (89.7\%), while the lowest
was recorded for avoiding shops (39.0\%), avoiding going out (49.3\%), and
wearing a face mask outside the home (52.5\%).

\subsection{Adherence trends across six-wave blocks}\label{adherence-trends-across-six-wave-blocks}

\begin{figure}[H]

{\centering \includegraphics[width=1\linewidth]{../graphs and tables/02_6week_cost_group_trend} 

}

\caption{Weighted adherence trends by behavioural-cost group across six-wave blocks.}\label{fig:cost-group-trends-figure}
\end{figure}

Figure 2 presents weighted adherence across the nine six-wave blocks.
Low-cost adherence ranged from 84.0\% to 87.0\%, medium-cost adherence
ranged from 55.4\% to 76.6\%, and high-cost adherence ranged from 41.0\% to
78.6\%. The largest adjacent increase for the medium- and high-cost
groups occurred between qweeks 31--36 and qweeks 37--42.

\subsubsection{Individual medium-cost behaviour trends}\label{individual-medium-cost-behaviour-trends}

\begin{figure}[H]

{\centering \includegraphics[width=1\linewidth]{../graphs and tables/04_medium_behaviour_trend} 

}

\caption{Weighted adherence trends for medium-cost protective behaviours across six-wave blocks.}\label{fig:medium-cost-trends-figure}
\end{figure}

Figure 3 presents weighted adherence for the five medium-cost
behaviours. Between qweeks 31--36 and qweeks 37--42, mask wearing
increased from 52.5\% to 78.3\%, avoiding guests from 55.8\% to 72.1\%,
avoiding crowded areas from 64.7\% to 76.6\%, avoiding public transport
from 61.1\% to 72.9\%, and avoiding people with COVID-like symptoms from
76.0\% to 83.0\%.

\subsubsection{Individual high-cost behaviour trends}\label{individual-high-cost-behaviour-trends}

\begin{figure}[H]

{\centering \includegraphics[width=1\linewidth]{../graphs and tables/04_high_behaviour_trend} 

}

\caption{Weighted adherence trends for high-cost protective behaviours across six-wave blocks.}\label{fig:high-cost-trends-figure}
\end{figure}

Figure 4 presents weighted adherence for the six high-cost behaviours.
Between qweeks 31--36 and qweeks 37--42, adherence increased from 49.5\% to
66.7\% for avoiding small gatherings, 55.3\% to 71.4\% for avoiding
medium-sized gatherings, 62.8\% to 77.5\% for avoiding large gatherings,
45.4\% to 57.4\% for avoiding going out, 36.4\% to 48.1\% for avoiding
shops, and 52.7\% to 61.9\% for avoiding healthcare settings.

\subsection{Behaviour-level change between selected periods}\label{behaviour-level-change-between-selected-periods}

\begin{longtable}[t]{llrrr}
\caption{\label{tab:behaviour-period-change-table}Change in weighted adherence between qweeks 31--36 and qweeks 37--42 by behaviour.}\\
\toprule
Behaviour & Cost group & Earlier (\%) & Later (\%) & Change (pp)\\
\midrule
Wore mask outside home & Medium & 52.5 & 78.3 & 25.8\\
Avoided small gatherings & High & 49.5 & 66.7 & 17.2\\
Avoided guests at home & Medium & 55.8 & 72.1 & 16.3\\
Avoided medium gatherings & High & 55.3 & 71.4 & 16.2\\
Avoided large gatherings & High & 62.8 & 77.5 & 14.7\\
\addlinespace
Avoided going out & High & 45.4 & 57.4 & 12.0\\
Avoided crowded areas & Medium & 64.7 & 76.6 & 11.9\\
Avoided public transport & Medium & 61.1 & 72.9 & 11.8\\
Avoided shops & High & 36.4 & 48.1 & 11.6\\
Avoided healthcare settings & High & 52.7 & 61.9 & 9.2\\
\addlinespace
Avoided symptomatic people & Medium & 76.0 & 83.0 & 7.0\\
Covered cough/sneeze & Low & 91.0 & 92.7 & 1.7\\
Washed hands & Low & 89.5 & 90.6 & 1.1\\
Used sanitiser & Low & 76.8 & 77.7 & 0.9\\
\bottomrule
\end{longtable}

Table \ref{tab:behaviour-period-change-table} reports the change
between qweeks 31--36 and qweeks 37--42 for all 14 behaviours. Mask
wearing recorded the largest increase at 25.8 percentage points,
followed by avoiding small gatherings at 17.2 percentage points,
avoiding guests at home at 16.3 percentage points, avoiding medium-sized
gatherings at 16.2 percentage points, and avoiding large gatherings at
14.7 percentage points. The three low-cost behaviours changed by between
0.9 and 1.7 percentage points.

\subsection{State-level descriptive changes}\label{state-level-descriptive-changes}

\begin{figure}[H]

{\centering \includegraphics[width=1\linewidth]{../graphs and tables/03_state_surge_change} 

}

\caption{Change in weighted adherence between qweeks 31--36 and qweeks 37--42 by state or territory for the medium- and high-cost groups.}\label{fig:state-period-change-figure}
\end{figure}

Figure 5 presents the change in medium- and high-cost adherence between
the selected periods across states and territories. Medium-cost changes
ranged from -16.2 to 29.9 percentage points, while high-cost changes
ranged from -14.4 to 32.4 percentage points. Low-cost changes are
reported in Appendix Table 12.

\subsection{Logistic regression results}\label{logistic-regression-results}

\begingroup
\footnotesize

\begin{longtable}[t]{lrrll}
\caption{\label{tab:model-3-period-results}Later-period coefficients and odds ratios for New South Wales from the state-by-period interaction models.}\\
\toprule
Behaviour & Period coefficient & Odds ratio & 95\% confidence interval & P-value\\
\midrule
Wearing a face mask outside the home & 2.2449 & 9.44 & 7.75--11.50 & <0.001\\
Avoiding guests at home & 1.5462 & 4.69 & 3.97--5.54 & <0.001\\
Avoiding small gatherings & 1.4214 & 4.14 & 3.56--4.83 & <0.001\\
Avoiding medium-sized gatherings & 1.4860 & 4.42 & 3.75--5.21 & <0.001\\
Avoiding large gatherings & 1.6567 & 5.24 & 4.35--6.32 & <0.001\\
\bottomrule
\end{longtable}

\endgroup

Table \ref{tab:model-3-period-results} reports the later-period
coefficients and odds ratios for New South Wales from the five
state-by-period interaction models. All five period coefficients were
positive and statistically significant at \(p < 0.001\). The corresponding
odds ratios ranged from 4.14 for avoiding small gatherings to 9.44 for
mask wearing. Complete state main effects and interaction coefficients
are reported in Appendix Tables 13--17.

\subsection{Model comparison results}\label{model-comparison-results}

\begin{longtable}[t]{lrrl}
\caption{\label{tab:model-comparison-table}Likelihood-ratio tests comparing the period-plus-state model (Model 2) with the state-by-period interaction model (Model 3).}\\
\toprule
Behaviour & LR statistic & Degrees of freedom & P-value\\
\midrule
Mask wearing & 662.42 & 7 & <0.001\\
Avoiding guests & 338.32 & 7 & <0.001\\
Avoiding small gatherings & 286.52 & 7 & <0.001\\
Avoiding medium gatherings & 338.39 & 7 & <0.001\\
Avoiding large gatherings & 345.28 & 7 & <0.001\\
\bottomrule
\end{longtable}

Table \ref{tab:model-comparison-table} reports the likelihood-ratio
tests comparing Model 2 with Model 3. The likelihood-ratio statistics
ranged from 286.52 to 662.42, with seven degrees of freedom for each
comparison. All five tests had \(p < 0.001\). For each behaviour, the
state-by-period interaction model therefore fitted the data better than
the period-plus-state model.

\section{Discussion}\label{discussion}

\subsection{Behavioural cost and adherence over time}\label{behavioural-cost-and-adherence-over-time}

The results provide general support for the idea that higher-cost
protective behaviours were less stable over time than low-cost
behaviours. Across the nine six-wave blocks, adherence to low-cost
hygiene behaviours remained relatively high, ranging from
approximately 84\% to 87\%. Meanwhile, the medium- and high-cost groups
showed substantially greater variation across the survey period.

This difference was also visible between qweeks 31--36 and qweeks
37--42. Covering coughs or sneezes, washing hands, and using hand
sanitiser increased by only 1.7, 1.1, and 0.9 percentage points,
respectively. These behaviours already had relatively high adherence and
could generally be performed without changing movement or
social contact. Their stability is consistent with earlier Australian
research showing that hygiene behaviours remained comparatively stable
while physical-distancing behaviours fluctuated over time (\citeproc{ref-ayre2021}{Ayre et al. 2021, 1}).

However, behavioural cost did not completely shape adherence. Some
medium- and high-cost behaviours increased sharply during the later
period. This indicates that the burden of a behaviour can be outweighed
by other influences, including outbreak severity, perceived infection
risk, social expectations, and public-health rules. This is consistent
with evidence that response costs do not affect all protective
behaviours in the same way (\citeproc{ref-hinssen2023}{Hinssen and Dohle 2023, sec. 4}). Behavioural cost
should therefore be understood as one influence on adherence than
as a complete explanation of behaviour.

\subsection{Behaviour-specific patterns and public-health context}\label{behaviour-specific-patterns-and-public-health-context}

Mask wearing showed the largest increase between the selected periods,
rising by 25.8 percentage points. The next largest changes were seen for
avoiding small gatherings at 17.2 percentage points, avoiding guests at
home at 16.3 percentage points, avoiding medium-sized gatherings at 16.2
percentage points, and avoiding large gatherings at 14.7 percentage
points.

Mask wearing differed from many other behaviours because it could be
directly required in particular settings. During the later period, mask
requirements were introduced or strengthened in New South Wales,
Victoria, and the Australian Capital Territory (\citeproc{ref-nsw_health_2021}{NSW Health 2021, paras. 1--3}; \citeproc{ref-vic_premier_2021}{State Government of Victoria 2021, paras. 1--7}; \citeproc{ref-act_gov_2021}{Barr and Stephen-Smith 2021, paras. 1--6}). This
provides a plausible explanation for why mask adherence increased more
than the other behaviours. Australian research has also shown that mask
mandates altered important determinants of adherence to protective
health behaviours (\citeproc{ref-ryan2025}{Ryan et al. 2025, 1}).

Avoiding guests and gatherings also increased substantially.
Restrictions introduced in several jurisdictions limited household
visitors, private and public gatherings, and permitted reasons for
leaving home (\citeproc{ref-vic_premier_2021}{State Government of Victoria 2021, paras. 1--7}; \citeproc{ref-act_gov_2021}{Barr and Stephen-Smith 2021, paras. 1--6}; \citeproc{ref-abc_sa_2021}{ABC News 2021, paras. 1--12}). These rules directly corresponded to
the behaviours measured in the survey, which may help explain their
simultaneous increases.

Changes in other movement-related behaviours were smaller. Avoiding
going out, shops, crowded areas, public transport, and healthcare
settings required greater changes to normal activities. However, some
movement remained necessary because restrictions continued to permit
essential shopping, healthcare, care responsibilities, and authorised
work (\citeproc{ref-vic_premier_2021}{State Government of Victoria 2021, paras. 1--7}; \citeproc{ref-act_gov_2021}{Barr and Stephen-Smith 2021, paras. 1--6}; \citeproc{ref-abc_sa_2021}{ABC News 2021, paras. 1--12}). This may explain why adherence to avoiding
shops and healthcare settings remained lower than adherence to avoiding
gatherings.

These interpretations remain cautious because the selected periods
contained several simultaneous changes. Outbreak conditions,
restrictions, media coverage, perceived risk, and social expectations
changed together. The analysis therefore cannot separate the effect of a
particular mandate or restriction from these other influences.

\subsection{State-level variation and model interpretation}\label{state-level-variation-and-model-interpretation}

The direction and magnitude of adherence changes varied across
Australian states. During the selected qweek period, medium-cost
adherence changes ranged from a decrease of 16.2 percentage points in
Western Australia to an increase of 29.9 percentage points in the
Australian Capital Territory. High-cost changes ranged from a decrease
of 14.4 percentage points in Western Australia to an increase of 32.4
percentage points in the Australian Capital Territory.

New South Wales and Victoria also recorded substantial increases in
medium- and high-cost adherence. These jurisdictions experienced major
Delta outbreaks and stronger restrictions during the later period. The
Australian Capital Territory entered lockdown in August 2021 and
introduced mandatory masks, retail restrictions, and limits on household
and social contact (\citeproc{ref-act_gov_2021}{Barr and Stephen-Smith 2021, paras. 1--6}). These conditions
correspond closely to the behaviours that increased.

Western Australia showed the opposite pattern, with decreases in both
medium- and high-cost adherence. During part of the selected period,
Western Australia was removing mask and crowd restrictions following an
absence of new local cases (\citeproc{ref-abc_wa_2021}{Perpitch 2021, paras. 1--5}). This contrast is
consistent with adherence changing differently in jurisdictions
experiencing different outbreak and policy conditions.

The regression results provided further evidence that the period change
differed across jurisdictions. For all five behaviours, the
state-by-period interaction model fitted significantly better than the
model containing only period and state main effects, with all
likelihood-ratio tests having \(p < 0.001\). This means that allowing each
jurisdiction to have a different period change improved the models.

For New South Wales, the estimated later-period odds ratios ranged from
4.14 for avoiding small gatherings to 9.44 for mask wearing. These
values represent multiplicative changes in the odds of adherence. For
example, the mask odds ratio of 9.44 means that the estimated odds of
reporting mask adherence in New South Wales were approximately 9.44
times higher during the later period than during the earlier period.

The interaction terms compared each state's period change with
the New South Wales change. A negative interaction coefficient indicated
that the increase was smaller than the New South Wales increase, while a
positive coefficient indicated a larger change. However, the models
identify statistical differences between jurisdictions and they also do not
show that particular policies caused those differences.

\subsection{Limitations}\label{limitations}

The Behaviour Tracker data were repeated cross-sectional. Different respondents were surveyed across waves, so the findings represent population-level differences between periods rather than changes within the same individuals.

The comparison between qweeks 31--36 and qweeks 37--42 was selected
after examining the descriptive trends. It was therefore exploratory and
should not be interpreted as a pre-specified evaluation of a particular
policy.

The timing of the changes in adherence overlapped with the Delta outbreak and changes
in public-health measures across different states, but the analysis does not establish
causation. Outbreak severity, perceived risk, media coverage, social expectations, and other unmeasured factors also have influenced adherence.

The behavioural-cost classification was developed specifically for this
project and involved literature based judgement. The burden associated with a behaviour
may differ between individuals according to their employment, health,
household circumstances, transport access, and financial position.

Adherence was measured through self-reported survey responses.
Respondents may have understood the questions differently or reported
socially desirable behaviour. Survey weighting improves the population
relevance of the estimates but cannot remove every form of sampling or
response bias.

Grouping the survey waves into six-wave blocks made the broad trends
easier to observe, but it may have concealed shorter changes within
individual waves. Estimates for smaller jurisdictions should also be
interpreted cautiously because they were based on fewer responses. The
state-level percentage-point comparisons were descriptive and did not
test whether each jurisdiction's individual change was statistically
significant. They also cannot attribute the national adherence pattern
causally to particular jurisdictions.

Finally, the regression analysis covered only the five behaviours with
the largest descriptive increases and did not adjust for demographic
characteristics. The supplied respondent-level weights were used for the
descriptive estimates and were implemented as frequency weights in the
regression models. The analytical dataset did not include the complete
survey-design information required to reconstruct the weights or conduct
a full design-based survey analysis. The resulting regression estimates
should therefore be interpreted as weighted model-based estimates rather
than design-based causal effects.

\subsection{Future research}\label{future-research}

Future research could include detailed state-level measures of outbreak
severity, policy timing, and restriction intensity. This would allow the
relationship between adherence and particular public-health measures to
be examined more directly.

Further models could adjust for demographic characteristics such as age,
gender and employment. Alternative behavioural-cost classifications could also be tested to determine whether the conclusions depend on the behavioural cost scoring framework.

Longitudinal data linking the same individuals across survey waves would
allow future studies to distinguish changes within individuals from
changes in the composition of the surveyed population.

\section{Conclusion}\label{conclusion}

Answering the research question, higher-cost COVID-safe
behaviours were generally less stable over time than low-cost hygiene
behaviours. Low-cost behaviours remained comparatively high and stable,
while medium- and high-cost behaviours showed greater variation.
However, behavioural cost alone did not explain every behaviour-specific
pattern.

The largest increases between qweeks 31--36 and qweeks 37--42 occurred
for mask wearing, avoiding guests, and avoiding small, medium-sized, and
large gatherings. The regression models also showed that the size of
these changes varied across states and territories.

The timing of these increases in adherence coincided with the 2021 Delta
outbreak and stronger public-health measures across several states and
territories. However, because the analysis used repeated cross-sectional
data and the periods were selected exploratorily, it cannot establish
that individual policies caused the observed changes. Overall, the
findings suggest that the adherence was related to both the burden of the
behaviour and the jurisdiction-specific outbreak and policy context.

\clearpage

\section{Acknowledgements}\label{acknowledgements}

I would like to acknowledge the guidance and feedback provided by my
project supervisor throughout the development of this project.

The analysis code, report source files, figures, tables, and numerical
outputs for this project are available in the following public GitHub
repository: {[}LINK{]}.

OpenAI's ChatGPT was used as a learning and editing aid to support my understanding of statistical concepts and to provide feedback on clarity, structure, grammar, and wording. I independently reviewed and revised the final report, checked the numerical results against the analysis code and saved outputs, and accept responsibility for the submitted work.

\section{References}\label{references}

\phantomsection\label{refs}
\begin{CSLReferences}{1}{0}
\bibitem[\citeproctext]{ref-abc_sa_2021}
ABC News. 2021. {``South Australia 'Moves into Lockdown' After Five COVID Cases Associated with Modbury Hospital Cluster.''} \url{https://www.abc.net.au/news/2021-07-20/south-australia-moves-into-lockdown/100306696}.

\bibitem[\citeproctext]{ref-abs_delta_2022}
Australian Bureau of Statistics. 2022. {``Effects of COVID-19 Strains on the Australian Economy.''} \url{https://www.abs.gov.au/articles/effects-covid-19-strains-australian-economy}.

\bibitem[\citeproctext]{ref-ayre2021}
Ayre, Julie, Erin Cvejic, Kirsten McCaffery, Tessa Copp, Samuel Cornell, Rachael H. Dodd, Kristen Pickles, et al. 2021. {``Contextualising COVID-19 Prevention Behaviour over Time in Australia: Patterns and Long-Term Predictors from April to July 2020 in an Online Social Media Sample.''} \emph{PLOS ONE} 16 (6): e0253930. \url{https://doi.org/10.1371/journal.pone.0253930}.

\bibitem[\citeproctext]{ref-act_gov_2021}
Barr, Andrew, and Rachel Stephen-Smith. 2021. {``COVID-19 Statement: ACT to Enter Seven Day Lockdown.''} ACT Government. \url{https://www.cmtedd.act.gov.au/open_government/inform/act_government_media_releases/rachel-stephen-smith-mla-media-releases/2021/covid-19-statement-act-to-enter-seven-day-lockdown}.

\bibitem[\citeproctext]{ref-hinssen2023}
Hinssen, Marina, and Simone Dohle. 2023. {``Personal Protective Behaviors in Response to COVID-19: A Longitudinal Application of Protection Motivation Theory.''} \emph{Frontiers in Psychology} 14: 1195607. \url{https://doi.org/10.3389/fpsyg.2023.1195607}.

\bibitem[\citeproctext]{ref-hosmer2013}
Hosmer, David W., Stanley Lemeshow, and Rodney X. Sturdivant. 2013. \emph{Applied Logistic Regression}. 3rd ed. Hoboken, NJ: John Wiley \& Sons.

\bibitem[\citeproctext]{ref-jones2020}
Jones, S. P., Imperial College London Big Data Analytical Unit, and YouGov Plc. 2020. {``Imperial College London YouGov COVID Data Hub.''} Version 1.0, YouGov Plc. \url{https://github.com/YouGov-Data/covid-19-tracker}.

\bibitem[\citeproctext]{ref-nsw_health_2021}
NSW Health. 2021. {``Fighting the Delta Outbreak with New Restrictions for Local Government Areas (LGAs) of Concern.''} \url{https://www.health.nsw.gov.au/news/Pages/20210730_01.aspx}.

\bibitem[\citeproctext]{ref-abc_wa_2021}
Perpitch, Nicolas. 2021. {``Perth Set to Return to Pre-Lockdown Life, Small Businesses Offered \$3,000 Support Grants.''} ABC News. \url{https://www.abc.net.au/news/2021-07-09/wa-returns-to-pre-lockdown-life/100280650}.

\bibitem[\citeproctext]{ref-ryan2025}
Ryan, Matthew, Jinjing Ye, Justin Sexton, Roslyn I. Hickson, and Emily Brindal. 2025. {``Face Mask Mandates Alter Major Determinants of Adherence to Protective Health Behaviours in Australia.''} \emph{Royal Society Open Science} 12: 241941. \url{https://doi.org/10.1098/rsos.241941}.

\bibitem[\citeproctext]{ref-seabold2010}
Seabold, Skipper, and Josef Perktold. 2010. {``Statsmodels: Econometric and Statistical Modeling with Python.''} In \emph{Proceedings of the 9th Python in Science Conference}, 92--96. \url{https://doi.org/10.25080/Majora-92bf1922-011}.

\bibitem[\citeproctext]{ref-vic_premier_2021}
State Government of Victoria. 2021. {``Seven Day Lockdown to Keep Victorians Safe.''} \url{https://www.premier.vic.gov.au/site-4/seven-day-lockdown-keep-victorians-safe}.

\bibitem[\citeproctext]{ref-virtanen2020}
Virtanen, Pauli, Ralf Gommers, Travis E. Oliphant, Matt Haberland, Tyler Reddy, David Cournapeau, Evgeni Burovski, et al. 2020. {``SciPy 1.0: Fundamental Algorithms for Scientific Computing in Python.''} \emph{Nature Methods} 17 (3): 261--72. \url{https://doi.org/10.1038/s41592-019-0686-2}.

\end{CSLReferences}

\clearpage

\section{Appendix}\label{appendix}

\subsection{Data preparation and code references}\label{data-preparation-and-code-references}

The initial data preparation was carried out in \texttt{Data\ cleaning.ipynb}.
The dataset structure and data-quality checks are contained under
``Initial Data Quality Check'' and ``Final dataset structure check''.
Selection and renaming of the analysis variables are documented under
``Define important columns for my analysis'' and ``Rename important columns
to usable names''. The creation of the 14 binary adherence outcomes is
documented under ``Create binary adherence variables for each behaviour''.

The cleaned wide-format dataset produced by these steps was saved as
\texttt{UGov\_Data.csv}. These are repository-relative file names; the complete
files are available through the public GitHub repository identified in
the Acknowledgements.

\small

\subsection{Behavioural-cost scoring matrix}\label{behavioural-cost-scoring-matrix}

Each behaviour was scored from 1 to 3 across practical effort, social
disruption, lifestyle restriction, emotional burden, and financial
burden. The five scores were summed to obtain the total burden score.
Total scores of 5--6 were classified as low cost, 7--9 as medium cost,
and 10--15 as high cost. Table \ref{tab:cost-scoring} presents the
complete scoring matrix used in the analysis.

\textbf{Code and source files:} The scoring rationale is documented in
\texttt{Classification\ of\ Protection\ Behaviours.docx}. The grouping used in the
analysis is implemented in \texttt{EDA.ipynb}, under ``Behaviour columns in wide
format''. The table below is generated from
\texttt{graphs\ and\ tables/appendix\_A1\_cost\_scoring.csv}.

\begin{landscape}

\begin{table}

\caption{\label{tab:cost-scoring}Behavioural-cost scoring matrix used to classify the 14 protective behaviours.}
\centering
\begin{tabular}[t]{lrrrrrrl}
\toprule
Behaviour & Practical & Social & Lifestyle & Emotional & Financial & Total & Group\\
\midrule
Covering nose/mouth when sneezing or coughing & 1 & 1 & 1 & 1 & 1 & 5 & Low\\
Washing hands with soap and water & 1 & 1 & 1 & 1 & 1 & 5 & Low\\
Using hand sanitiser & 1 & 1 & 1 & 1 & 2 & 6 & Low\\
Wore mask outside home & 2 & 2 & 1 & 2 & 1 & 8 & Medium\\
Avoiding contact with symptomatic people & 1 & 2 & 1 & 2 & 1 & 7 & Medium\\
\addlinespace
Avoiding crowded areas & 1 & 2 & 2 & 2 & 1 & 8 & Medium\\
Avoiding public transport & 2 & 1 & 2 & 2 & 1 & 8 & Medium\\
Avoiding having guests at home & 1 & 2 & 2 & 2 & 1 & 8 & Medium\\
Avoiding small social gatherings & 2 & 2 & 2 & 3 & 1 & 10 & High\\
Avoiding medium social gatherings & 1 & 3 & 2 & 3 & 1 & 10 & High\\
\addlinespace
Avoiding large social gatherings & 1 & 3 & 2 & 3 & 1 & 10 & High\\
Avoiding going out in general & 2 & 3 & 3 & 3 & 1 & 12 & High\\
Avoiding going to shops & 2 & 2 & 3 & 2 & 1 & 10 & High\\
Avoiding going to healthcare settings & 2 & 1 & 3 & 3 & 1 & 10 & High\\
\bottomrule
\end{tabular}
\end{table}

\end{landscape}

\subsection{Survey-wave and date mapping}\label{survey-wave-and-date-mapping}

Table \ref{tab:qweek-dates} reports the observed response-date range,
six-wave block, and unweighted number of responses for each \texttt{qweek}.

\textbf{Code and saved output:} The numeric survey-wave ordering was created
in \texttt{Data\ cleaning.ipynb} during the construction of \texttt{week\_num}. The
observed date ranges were obtained from \texttt{response\_time}. The table below
is generated from \texttt{graphs\ and\ tables/appendix\_A2\_qweek\_dates.csv}.

\begin{longtable}[t]{rlllr}
\caption{\label{tab:qweek-dates}Observed response-date range and six-wave block for each survey wave.}\\
\toprule
qweek & Survey start date & Survey end date & Six-wave block & Unweighted n\\
\midrule
1 & 01 Apr 2020 & 03 Apr 2020 & 1--6 & 973\\
2 & 07 Apr 2020 & 09 Apr 2020 & 1--6 & 1006\\
3 & 15 Apr 2020 & 20 Apr 2020 & 1--6 & 1002\\
4 & 21 Apr 2020 & 24 Apr 2020 & 1--6 & 1004\\
5 & 30 Apr 2020 & 04 May 2020 & 1--6 & 1007\\
\addlinespace
6 & 06 May 2020 & 07 May 2020 & 1--6 & 1014\\
7 & 15 May 2020 & 18 May 2020 & 7--12 & 1008\\
8 & 03 Jun 2020 & 05 Jun 2020 & 7--12 & 1004\\
9 & 09 Jun 2020 & 11 Jun 2020 & 7--12 & 1005\\
10 & 24 Jun 2020 & 28 Jun 2020 & 7--12 & 1009\\
\addlinespace
11 & 08 Jul 2020 & 13 Jul 2020 & 7--12 & 1010\\
12 & 22 Jul 2020 & 27 Jul 2020 & 7--12 & 1001\\
13 & 06 Aug 2020 & 10 Aug 2020 & 13--18 & 1003\\
14 & 20 Aug 2020 & 25 Aug 2020 & 13--18 & 1013\\
15 & 02 Sep 2020 & 07 Sep 2020 & 13--18 & 1009\\
\addlinespace
16 & 16 Sep 2020 & 26 Sep 2020 & 13--18 & 1001\\
17 & 02 Oct 2020 & 12 Oct 2020 & 13--18 & 1000\\
18 & 14 Oct 2020 & 19 Oct 2020 & 13--18 & 1006\\
19 & 28 Oct 2020 & 04 Nov 2020 & 19--24 & 999\\
20 & 11 Nov 2020 & 17 Nov 2020 & 19--24 & 1001\\
\addlinespace
21 & 16 Dec 2020 & 30 Dec 2020 & 19--24 & 1002\\
22 & 30 Dec 2020 & 05 Jan 2021 & 19--24 & 924\\
23 & 13 Jan 2021 & 26 Jan 2021 & 19--24 & 1001\\
24 & 27 Jan 2021 & 08 Feb 2021 & 19--24 & 1003\\
25 & 10 Feb 2021 & 23 Feb 2021 & 25--30 & 1011\\
\addlinespace
26 & 24 Feb 2021 & 05 Mar 2021 & 25--30 & 966\\
27 & 10 Mar 2021 & 19 Mar 2021 & 25--30 & 944\\
28 & 24 Mar 2021 & 06 Apr 2021 & 25--30 & 1004\\
29 & 07 Apr 2021 & 19 Apr 2021 & 25--30 & 1004\\
30 & 21 Apr 2021 & 05 May 2021 & 25--30 & 1003\\
\addlinespace
31 & 05 May 2021 & 18 May 2021 & 31--36 & 969\\
32 & 19 May 2021 & 01 Jun 2021 & 31--36 & 1001\\
33 & 02 Jun 2021 & 14 Jun 2021 & 31--36 & 1002\\
34 & 16 Jun 2021 & 28 Jun 2021 & 31--36 & 1009\\
35 & 30 Jun 2021 & 13 Jul 2021 & 31--36 & 1006\\
\addlinespace
36 & 14 Jul 2021 & 27 Jul 2021 & 31--36 & 1008\\
37 & 28 Jul 2021 & 05 Aug 2021 & 37--42 & 1001\\
38 & 11 Aug 2021 & 17 Aug 2021 & 37--42 & 1001\\
39 & 23 Aug 2021 & 06 Sep 2021 & 37--42 & 1005\\
40 & 07 Sep 2021 & 17 Sep 2021 & 37--42 & 1011\\
\addlinespace
41 & 20 Sep 2021 & 03 Oct 2021 & 37--42 & 1006\\
42 & 04 Oct 2021 & 17 Oct 2021 & 37--42 & 1006\\
43 & 18 Oct 2021 & 26 Oct 2021 & 43--48 & 1007\\
44 & 01 Nov 2021 & 09 Nov 2021 & 43--48 & 1002\\
45 & 16 Nov 2021 & 23 Nov 2021 & 43--48 & 1021\\
\addlinespace
46 & 29 Nov 2021 & 10 Dec 2021 & 43--48 & 1006\\
47 & 14 Dec 2021 & 29 Dec 2021 & 43--48 & 978\\
48 & 27 Dec 2021 & 05 Jan 2022 & 43--48 & 1011\\
49 & 10 Jan 2022 & 24 Jan 2022 & 49--54 & 1010\\
50 & 24 Jan 2022 & 01 Feb 2022 & 49--54 & 1007\\
\addlinespace
51 & 07 Feb 2022 & 16 Feb 2022 & 49--54 & 897\\
52 & 21 Feb 2022 & 01 Mar 2022 & 49--54 & 930\\
53 & 07 Mar 2022 & 16 Mar 2022 & 49--54 & 996\\
54 & 21 Mar 2022 & 28 Mar 2022 & 49--54 & 1006\\
\bottomrule
\end{longtable}

\subsection{Complete descriptive estimates}\label{complete-descriptive-estimates}

\subsubsection{Overall adherence estimates}\label{overall-adherence-estimates}

\textbf{Code and saved output:} These estimates were generated in
\texttt{EDA.ipynb}, under ``Overall behaviour summary and plot''. The table is
generated from \texttt{graphs\ and\ tables/appendix\_A3\_overall\_adherence.csv}.

\begin{longtable}[t]{llrr}
\caption{\label{tab:overall-adherence-appendix}Overall weighted adherence estimates for the 14 protective behaviours.}\\
\toprule
Behaviour & Cost group & Valid responses & Weighted adherence (\%)\\
\midrule
Covered coughs or sneezes & Low & 53833 & 90.9\\
Washed hands & Low & 52827 & 89.7\\
Avoided symptomatic people & Medium & 53833 & 77.7\\
Used hand sanitiser & Low & 53833 & 76.6\\
Avoided crowded areas & Medium & 53833 & 70.8\\
\addlinespace
Avoided large gatherings & High & 53833 & 69.3\\
Avoided public transport & Medium & 53833 & 66.0\\
Avoided medium gatherings & High & 53833 & 60.4\\
Avoided guests at home & Medium & 52827 & 60.3\\
Avoided healthcare settings & High & 53833 & 54.9\\
\addlinespace
Avoided small gatherings & High & 53833 & 53.9\\
Wore mask outside home & Medium & 53833 & 52.5\\
Avoided going out & High & 53833 & 49.3\\
Avoided shops & High & 53833 & 39.0\\
\bottomrule
\end{longtable}

\subsubsection{Behavioural-cost group trends}\label{behavioural-cost-group-trends}

\textbf{Code and saved output:} These estimates were generated in
\texttt{EDA.ipynb}, under ``Create equal six-wave blocks''. The table is
generated from \texttt{graphs\ and\ tables/appendix\_A4\_cost\_group\_trends.csv}.

\begin{longtable}[t]{lrrr}
\caption{\label{tab:cost-trends-appendix}Weighted adherence estimates by behavioural-cost group and six-wave block.}\\
\toprule
Six-wave block & Low cost (\%) & Medium cost (\%) & High cost (\%)\\
\midrule
1--6 & 85.7 & 73.3 & 78.6\\
7--12 & 86.1 & 63.1 & 58.6\\
13--18 & 86.7 & 63.8 & 54.0\\
19--24 & 86.7 & 60.6 & 46.6\\
25--30 & 85.0 & 55.4 & 41.0\\
\addlinespace
31--36 & 85.8 & 62.0 & 50.4\\
37--42 & 87.0 & 76.6 & 63.8\\
43--48 & 84.0 & 65.6 & 47.7\\
49--54 & 84.0 & 68.7 & 49.3\\
\bottomrule
\end{longtable}

\subsubsection{Medium-cost behaviour trends}\label{medium-cost-behaviour-trends}

\textbf{Code and saved output:} These estimates were generated in
\texttt{EDA.ipynb}, under ``Individual behaviour trends by six-wave block''. The
table is generated from
\texttt{graphs\ and\ tables/appendix\_A5\_medium\_behaviour\_trends.csv}.

\begin{longtable}[t]{lllr}
\caption{\label{tab:medium-trends-appendix}Weighted adherence estimates for medium-cost behaviours by six-wave block.}\\
\toprule
Six-wave block & Behaviour & Cost group & Weighted adherence (\%)\\
\midrule
1--6 & Wore mask outside home & Medium & 24.6\\
1--6 & Avoided symptomatic people & Medium & 85.2\\
1--6 & Avoided crowded areas & Medium & 88.3\\
1--6 & Avoided public transport & Medium & 82.6\\
1--6 & Avoided guests at home & Medium & 85.7\\
\addlinespace
7--12 & Wore mask outside home & Medium & 24.3\\
7--12 & Avoided symptomatic people & Medium & 78.1\\
7--12 & Avoided crowded areas & Medium & 77.6\\
7--12 & Avoided public transport & Medium & 71.4\\
7--12 & Avoided guests at home & Medium & 64.0\\
\addlinespace
13--18 & Wore mask outside home & Medium & 47.7\\
13--18 & Avoided symptomatic people & Medium & 74.1\\
13--18 & Avoided crowded areas & Medium & 72.7\\
13--18 & Avoided public transport & Medium & 65.4\\
13--18 & Avoided guests at home & Medium & 59.2\\
\addlinespace
19--24 & Wore mask outside home & Medium & 48.3\\
19--24 & Avoided symptomatic people & Medium & 74.7\\
19--24 & Avoided crowded areas & Medium & 66.8\\
19--24 & Avoided public transport & Medium & 61.8\\
19--24 & Avoided guests at home & Medium & 51.2\\
\addlinespace
25--30 & Wore mask outside home & Medium & 43.5\\
25--30 & Avoided symptomatic people & Medium & 72.2\\
25--30 & Avoided crowded areas & Medium & 59.0\\
25--30 & Avoided public transport & Medium & 55.1\\
25--30 & Avoided guests at home & Medium & 47.0\\
\addlinespace
31--36 & Wore mask outside home & Medium & 52.5\\
31--36 & Avoided symptomatic people & Medium & 76.0\\
31--36 & Avoided crowded areas & Medium & 64.7\\
31--36 & Avoided public transport & Medium & 61.1\\
31--36 & Avoided guests at home & Medium & 55.8\\
\addlinespace
37--42 & Wore mask outside home & Medium & 78.3\\
37--42 & Avoided symptomatic people & Medium & 83.0\\
37--42 & Avoided crowded areas & Medium & 76.6\\
37--42 & Avoided public transport & Medium & 72.9\\
37--42 & Avoided guests at home & Medium & 72.1\\
\addlinespace
43--48 & Wore mask outside home & Medium & 73.0\\
43--48 & Avoided symptomatic people & Medium & 76.5\\
43--48 & Avoided crowded areas & Medium & 64.6\\
43--48 & Avoided public transport & Medium & 61.5\\
43--48 & Avoided guests at home & Medium & 52.4\\
\addlinespace
49--54 & Wore mask outside home & Medium & 80.6\\
49--54 & Avoided symptomatic people & Medium & 79.1\\
49--54 & Avoided crowded areas & Medium & 66.8\\
49--54 & Avoided public transport & Medium & 61.7\\
49--54 & Avoided guests at home & Medium & 53.9\\
\bottomrule
\end{longtable}

\subsubsection{High-cost behaviour trends}\label{high-cost-behaviour-trends}

\textbf{Code and saved output:} These estimates were generated in
\texttt{EDA.ipynb}, under ``Individual behaviour trends by six-wave block''. The
table is generated from
\texttt{graphs\ and\ tables/appendix\_A6\_high\_behaviour\_trends.csv}.

\begin{longtable}[t]{lllr}
\caption{\label{tab:high-trends-appendix}Weighted adherence estimates for high-cost behaviours by six-wave block.}\\
\toprule
Six-wave block & Behaviour & Cost group & Weighted adherence (\%)\\
\midrule
1--6 & Avoided small gatherings & High & 84.0\\
1--6 & Avoided medium gatherings & High & 89.7\\
1--6 & Avoided large gatherings & High & 91.9\\
1--6 & Avoided going out & High & 75.7\\
1--6 & Avoided shops & High & 58.8\\
\addlinespace
1--6 & Avoided healthcare settings & High & 71.7\\
7--12 & Avoided small gatherings & High & 56.6\\
7--12 & Avoided medium gatherings & High & 66.5\\
7--12 & Avoided large gatherings & High & 77.7\\
7--12 & Avoided going out & High & 53.5\\
\addlinespace
7--12 & Avoided shops & High & 42.4\\
7--12 & Avoided healthcare settings & High & 54.9\\
13--18 & Avoided small gatherings & High & 52.8\\
13--18 & Avoided medium gatherings & High & 60.8\\
13--18 & Avoided large gatherings & High & 71.7\\
\addlinespace
13--18 & Avoided going out & High & 49.7\\
13--18 & Avoided shops & High & 38.4\\
13--18 & Avoided healthcare settings & High & 50.6\\
19--24 & Avoided small gatherings & High & 44.0\\
19--24 & Avoided medium gatherings & High & 51.4\\
\addlinespace
19--24 & Avoided large gatherings & High & 62.8\\
19--24 & Avoided going out & High & 40.5\\
19--24 & Avoided shops & High & 31.1\\
19--24 & Avoided healthcare settings & High & 49.8\\
25--30 & Avoided small gatherings & High & 38.9\\
\addlinespace
25--30 & Avoided medium gatherings & High & 44.6\\
25--30 & Avoided large gatherings & High & 55.0\\
25--30 & Avoided going out & High & 34.3\\
25--30 & Avoided shops & High & 26.7\\
25--30 & Avoided healthcare settings & High & 46.3\\
\addlinespace
31--36 & Avoided small gatherings & High & 49.5\\
31--36 & Avoided medium gatherings & High & 55.3\\
31--36 & Avoided large gatherings & High & 62.8\\
31--36 & Avoided going out & High & 45.4\\
31--36 & Avoided shops & High & 36.4\\
\addlinespace
31--36 & Avoided healthcare settings & High & 52.7\\
37--42 & Avoided small gatherings & High & 66.7\\
37--42 & Avoided medium gatherings & High & 71.4\\
37--42 & Avoided large gatherings & High & 77.5\\
37--42 & Avoided going out & High & 57.4\\
\addlinespace
37--42 & Avoided shops & High & 48.1\\
37--42 & Avoided healthcare settings & High & 61.9\\
43--48 & Avoided small gatherings & High & 45.7\\
43--48 & Avoided medium gatherings & High & 51.4\\
43--48 & Avoided large gatherings & High & 60.4\\
\addlinespace
43--48 & Avoided going out & High & 42.4\\
43--48 & Avoided shops & High & 34.4\\
43--48 & Avoided healthcare settings & High & 51.9\\
49--54 & Avoided small gatherings & High & 46.7\\
49--54 & Avoided medium gatherings & High & 52.4\\
\addlinespace
49--54 & Avoided large gatherings & High & 63.0\\
49--54 & Avoided going out & High & 44.6\\
49--54 & Avoided shops & High & 34.4\\
49--54 & Avoided healthcare settings & High & 54.5\\
\bottomrule
\end{longtable}

\clearpage

\subsubsection{State-level estimates and sample sizes}\label{state-level-estimates-and-sample-sizes}

\textbf{Code and saved outputs:} The state-level estimates and changes were
generated in \texttt{EDA.ipynb}, under ``State-level changes between the
selected periods''. The tables are generated from
\texttt{graphs\ and\ tables/appendix\_A7a\_state\_period\_estimates.csv} and
\texttt{graphs\ and\ tables/appendix\_A7b\_state\_changes.csv}.

\begin{landscape}
\scriptsize


\begin{longtable}[t]{llrrrr}
\caption{\label{tab:state-estimates}Unweighted sample sizes and weighted cost-group adherence estimates by jurisdiction and period.}\\
\toprule
State or territory & Period & Unweighted n & Low cost (\%) & Medium cost (\%) & High cost (\%)\\
\midrule
Australian Capital Territory & qweeks 31--36 & 87 & 83.1 & 54.6 & 39.0\\
Australian Capital Territory & qweeks 37--42 & 99 & 86.3 & 84.5 & 71.4\\
New South Wales & qweeks 31--36 & 1802 & 87.5 & 66.6 & 56.0\\
New South Wales & qweeks 37--42 & 1864 & 89.5 & 89.3 & 79.3\\
Northern Territory & qweeks 31--36 & 46 & 79.5 & 47.6 & 31.2\\
\addlinespace
Northern Territory & qweeks 37--42 & 41 & 83.6 & 54.7 & 44.7\\
Queensland & qweeks 31--36 & 1197 & 83.2 & 53.9 & 43.9\\
Queensland & qweeks 37--42 & 1187 & 85.4 & 66.9 & 49.5\\
South Australia & qweeks 31--36 & 614 & 84.3 & 50.3 & 41.1\\
South Australia & qweeks 37--42 & 596 & 85.9 & 71.5 & 48.6\\
\addlinespace
Tasmania & qweeks 31--36 & 121 & 81.1 & 37.3 & 27.9\\
Tasmania & qweeks 37--42 & 123 & 81.8 & 42.1 & 31.9\\
Victoria & qweeks 31--36 & 1535 & 87.0 & 70.6 & 55.6\\
Victoria & qweeks 37--42 & 1555 & 88.9 & 86.8 & 75.1\\
Western Australia & qweeks 31--36 & 593 & 85.1 & 57.6 & 46.8\\
\addlinespace
Western Australia & qweeks 37--42 & 565 & 80.0 & 41.4 & 32.4\\
\bottomrule
\end{longtable}


\begin{longtable}[t]{lrrr}
\caption{\label{tab:state-changes}Percentage-point change in weighted cost-group adherence between the two selected periods.}\\
\toprule
State or territory & Low cost change (pp) & Medium cost change (pp) & High cost change (pp)\\
\midrule
Australian Capital Territory & 3.2 & 29.9 & 32.4\\
New South Wales & 2.0 & 22.7 & 23.3\\
Northern Territory & 4.1 & 7.1 & 13.5\\
Queensland & 2.2 & 13.0 & 5.6\\
South Australia & 1.6 & 21.2 & 7.5\\
\addlinespace
Tasmania & 0.7 & 4.8 & 4.0\\
Victoria & 1.9 & 16.2 & 19.5\\
Western Australia & -5.1 & -16.2 & -14.4\\
\bottomrule
\end{longtable}

\small
\end{landscape}
\clearpage

\subsection{Complete Model 3 coefficient outputs}\label{complete-model-3-coefficient-outputs}

\textbf{Code and saved outputs:} The five models were fitted in
\texttt{Modelling.ipynb}, within the behaviour-specific Model 3 sections.
Models were implemented in Python using \texttt{statsmodels.formula.api.glm()}
with a binomial family and logit link. Respondent weights were passed
through the \texttt{freq\_weights} argument, and HC1 robust covariance estimates
were used when calculating standard errors, confidence intervals, and
coefficient p-values. The tables below are generated from
\texttt{graphs\ and\ tables/appendix\_A8\_model3\_mask\_wearing.csv} through
\texttt{graphs\ and\ tables/appendix\_A12\_model3\_avoiding\_large\_gatherings.csv}.

New South Wales and qweeks 31--36 were the reference categories. State
main effects compare each jurisdiction with New South Wales during
qweeks 31--36. Interaction terms show how each jurisdiction's period
change differed from the New South Wales period change.

\subsubsection{Wearing a face mask outside the home}\label{wearing-a-face-mask-outside-the-home}

\begin{longtable}[t]{lrrll}
\caption{\label{tab:model3-mask-appendix}Complete Model 3 coefficient output for mask wearing.}\\
\toprule
Term & Coefficient & Odds ratio & 95\% CI & P-value\\
\midrule
Intercept (NSW, qweeks 31--36) & 0.338 & 1.402 & 1.281--1.536 & <0.001\\
ACT vs NSW (qweeks 31--36) & -0.592 & 0.553 & 0.370--0.828 & 0.004\\
NT vs NSW (qweeks 31--36) & -1.672 & 0.188 & 0.098--0.358 & <0.001\\
QLD vs NSW (qweeks 31--36) & -0.935 & 0.393 & 0.338--0.456 & <0.001\\
SA vs NSW (qweeks 31--36) & -1.419 & 0.242 & 0.191--0.307 & <0.001\\
\addlinespace
TAS vs NSW (qweeks 31--36) & -2.850 & 0.058 & 0.030--0.112 & <0.001\\
VIC vs NSW (qweeks 31--36) & 0.709 & 2.032 & 1.758--2.349 & <0.001\\
WA vs NSW (qweeks 31--36) & -0.529 & 0.589 & 0.491--0.707 & <0.001\\
Later period: qweeks 37--42 (NSW) & 2.245 & 9.439 & 7.754--11.491 & <0.001\\
Later period x ACT & -0.703 & 0.495 & 0.260--0.944 & 0.033\\
\addlinespace
Later period x NT & -1.427 & 0.240 & 0.102--0.565 & 0.001\\
Later period x QLD & -0.677 & 0.508 & 0.391--0.661 & <0.001\\
Later period x SA & 0.927 & 2.526 & 1.655--3.856 & <0.001\\
Later period x TAS & -2.202 & 0.111 & 0.043--0.283 & <0.001\\
Later period x VIC & -0.741 & 0.477 & 0.354--0.641 & <0.001\\
\addlinespace
Later period x WA & -3.528 & 0.029 & 0.021--0.041 & <0.001\\
\bottomrule
\end{longtable}

\clearpage

\subsubsection{Avoiding guests at home}\label{avoiding-guests-at-home}

\begin{longtable}[t]{lrrll}
\caption{\label{tab:model3-guests-appendix}Complete Model 3 coefficient output for avoiding guests at home.}\\
\toprule
Term & Coefficient & Odds ratio & 95\% CI & P-value\\
\midrule
Intercept (NSW, qweeks 31--36) & 0.466 & 1.594 & 1.454--1.748 & <0.001\\
ACT vs NSW (qweeks 31--36) & -0.511 & 0.600 & 0.402--0.896 & 0.012\\
NT vs NSW (qweeks 31--36) & -0.393 & 0.675 & 0.398--1.145 & 0.145\\
QLD vs NSW (qweeks 31--36) & -0.611 & 0.543 & 0.469--0.628 & <0.001\\
SA vs NSW (qweeks 31--36) & -0.711 & 0.491 & 0.397--0.608 & <0.001\\
\addlinespace
TAS vs NSW (qweeks 31--36) & -1.340 & 0.262 & 0.177--0.387 & <0.001\\
VIC vs NSW (qweeks 31--36) & 0.084 & 1.087 & 0.947--1.248 & 0.234\\
WA vs NSW (qweeks 31--36) & -0.388 & 0.679 & 0.565--0.815 & <0.001\\
Later period: qweeks 37--42 (NSW) & 1.546 & 4.693 & 3.975--5.541 & <0.001\\
Later period x ACT & -0.298 & 0.742 & 0.397--1.389 & 0.351\\
\addlinespace
Later period x NT & -1.697 & 0.183 & 0.086--0.389 & <0.001\\
Later period x QLD & -1.192 & 0.304 & 0.241--0.383 & <0.001\\
Later period x SA & -1.135 & 0.322 & 0.234--0.442 & <0.001\\
Later period x TAS & -1.198 & 0.302 & 0.175--0.521 & <0.001\\
Later period x VIC & -0.132 & 0.877 & 0.685--1.122 & 0.295\\
\addlinespace
Later period x WA & -2.203 & 0.110 & 0.083--0.146 & <0.001\\
\bottomrule
\end{longtable}

\clearpage

\subsubsection{Avoiding small gatherings}\label{avoiding-small-gatherings}

\begin{longtable}[t]{lrrll}
\caption{\label{tab:model3-small-appendix}Complete Model 3 coefficient output for avoiding small gatherings.}\\
\toprule
Term & Coefficient & Odds ratio & 95\% CI & P-value\\
\midrule
Intercept (NSW, qweeks 31--36) & 0.260 & 1.296 & 1.185--1.419 & <0.001\\
ACT vs NSW (qweeks 31--36) & -0.686 & 0.504 & 0.335--0.758 & 0.001\\
NT vs NSW (qweeks 31--36) & -0.811 & 0.444 & 0.257--0.768 & 0.004\\
QLD vs NSW (qweeks 31--36) & -0.622 & 0.537 & 0.463--0.621 & <0.001\\
SA vs NSW (qweeks 31--36) & -0.756 & 0.470 & 0.378--0.583 & <0.001\\
\addlinespace
TAS vs NSW (qweeks 31--36) & -1.457 & 0.233 & 0.153--0.354 & <0.001\\
VIC vs NSW (qweeks 31--36) & -0.007 & 0.993 & 0.867--1.136 & 0.913\\
WA vs NSW (qweeks 31--36) & -0.509 & 0.601 & 0.501--0.722 & <0.001\\
Later period: qweeks 37--42 (NSW) & 1.421 & 4.143 & 3.558--4.824 & <0.001\\
Later period x ACT & -0.107 & 0.899 & 0.491--1.645 & 0.729\\
\addlinespace
Later period x NT & -1.090 & 0.336 & 0.156--0.723 & 0.005\\
Later period x QLD & -1.091 & 0.336 & 0.269--0.420 & <0.001\\
Later period x SA & -1.095 & 0.335 & 0.245--0.458 & <0.001\\
Later period x TAS & -1.305 & 0.271 & 0.150--0.489 & <0.001\\
Later period x VIC & -0.147 & 0.863 & 0.691--1.079 & 0.197\\
\addlinespace
Later period x WA & -1.981 & 0.138 & 0.104--0.182 & <0.001\\
\bottomrule
\end{longtable}

\clearpage

\subsubsection{Avoiding medium-sized gatherings}\label{avoiding-medium-sized-gatherings}

\begin{longtable}[t]{lrrll}
\caption{\label{tab:model3-medium-appendix}Complete Model 3 coefficient output for avoiding medium-sized gatherings.}\\
\toprule
Term & Coefficient & Odds ratio & 95\% CI & P-value\\
\midrule
Intercept (NSW, qweeks 31--36) & 0.500 & 1.649 & 1.504--1.808 & <0.001\\
ACT vs NSW (qweeks 31--36) & -0.727 & 0.484 & 0.323--0.724 & <0.001\\
NT vs NSW (qweeks 31--36) & -1.052 & 0.349 & 0.202--0.604 & <0.001\\
QLD vs NSW (qweeks 31--36) & -0.576 & 0.562 & 0.486--0.651 & <0.001\\
SA vs NSW (qweeks 31--36) & -0.702 & 0.495 & 0.401--0.613 & <0.001\\
\addlinespace
TAS vs NSW (qweeks 31--36) & -1.226 & 0.294 & 0.201--0.429 & <0.001\\
VIC vs NSW (qweeks 31--36) & -0.076 & 0.927 & 0.808--1.063 & 0.279\\
WA vs NSW (qweeks 31--36) & -0.529 & 0.589 & 0.491--0.707 & <0.001\\
Later period: qweeks 37--42 (NSW) & 1.486 & 4.419 & 3.747--5.213 & <0.001\\
Later period x ACT & 0.235 & 1.265 & 0.655--2.445 & 0.484\\
\addlinespace
Later period x NT & -1.078 & 0.340 & 0.158--0.732 & 0.006\\
Later period x QLD & -1.234 & 0.291 & 0.231--0.367 & <0.001\\
Later period x SA & -1.202 & 0.301 & 0.219--0.412 & <0.001\\
Later period x TAS & -1.194 & 0.303 & 0.177--0.518 & <0.001\\
Later period x VIC & -0.053 & 0.948 & 0.745--1.208 & 0.667\\
\addlinespace
Later period x WA & -2.075 & 0.126 & 0.095--0.166 & <0.001\\
\bottomrule
\end{longtable}

\clearpage

\subsubsection{Avoiding large gatherings}\label{avoiding-large-gatherings}

\begin{longtable}[t]{lrrll}
\caption{\label{tab:model3-large-appendix}Complete Model 3 coefficient output for avoiding large gatherings.}\\
\toprule
Term & Coefficient & Odds ratio & 95\% CI & P-value\\
\midrule
Intercept (NSW, qweeks 31--36) & 0.721 & 2.057 & 1.870--2.263 & <0.001\\
ACT vs NSW (qweeks 31--36) & -0.613 & 0.542 & 0.362--0.810 & 0.003\\
NT vs NSW (qweeks 31--36) & -0.978 & 0.376 & 0.221--0.640 & <0.001\\
QLD vs NSW (qweeks 31--36) & -0.431 & 0.650 & 0.560--0.755 & <0.001\\
SA vs NSW (qweeks 31--36) & -0.549 & 0.577 & 0.466--0.715 & <0.001\\
\addlinespace
TAS vs NSW (qweeks 31--36) & -0.923 & 0.397 & 0.277--0.570 & <0.001\\
VIC vs NSW (qweeks 31--36) & 0.049 & 1.050 & 0.910--1.212 & 0.501\\
WA vs NSW (qweeks 31--36) & -0.363 & 0.696 & 0.577--0.838 & <0.001\\
Later period: qweeks 37--42 (NSW) & 1.657 & 5.242 & 4.351--6.315 & <0.001\\
Later period x ACT & 0.478 & 1.613 & 0.732--3.551 & 0.235\\
\addlinespace
Later period x NT & -1.178 & 0.308 & 0.144--0.660 & 0.002\\
Later period x QLD & -1.352 & 0.259 & 0.202--0.332 & <0.001\\
Later period x SA & -1.245 & 0.288 & 0.207--0.402 & <0.001\\
Later period x TAS & -1.532 & 0.216 & 0.128--0.365 & <0.001\\
Later period x VIC & -0.267 & 0.766 & 0.585--1.002 & 0.052\\
\addlinespace
Later period x WA & -2.276 & 0.103 & 0.077--0.138 & <0.001\\
\bottomrule
\end{longtable}

\clearpage

\subsection{Likelihood-ratio test output}\label{likelihood-ratio-test-output}

Table \ref{tab:lr-tests-appendix} reports the likelihood-ratio tests
comparing Model 2 with Model 3 for the five modelled behaviours.

\textbf{Code and saved output:} The model comparisons were performed in
\texttt{Modelling.ipynb}, under ``Final model-comparison section''.
Likelihood-ratio p-values were calculated using \texttt{scipy.stats.chi2.sf()}.
The table is generated from
\texttt{graphs\ and\ tables/appendix\_A13\_likelihood\_ratio\_tests.csv}.

\begin{longtable}[t]{lrrl}
\caption{\label{tab:lr-tests-appendix}Likelihood-ratio tests comparing Model 2 with Model 3.}\\
\toprule
Behaviour & LR statistic & Degrees of freedom & P-value\\
\midrule
Mask wearing & 662.42 & 7 & <0.001\\
Avoiding guests & 338.32 & 7 & <0.001\\
Avoiding small gatherings & 286.52 & 7 & <0.001\\
Avoiding medium gatherings & 338.39 & 7 & <0.001\\
Avoiding large gatherings & 345.28 & 7 & <0.001\\
\bottomrule
\end{longtable}

\clearpage

The analysis code, report source file, cleaned dataset, figures, tables,
and saved numerical outputs are available in the public GitHub
repository provided in the Acknowledgements section.

\normalsize

\end{document}
